% Section: Glossary, Part C: identities (introduction) | Part: backmatter | Depends on: all parts | Status: done | Original lines: none (added by rewriter)
%!TEX root = ../main.tex
\phantomsection\subsection*{Part C. Identities, with proofs}\addcontentsline{toc}{subsection}{Part C. Identities, with proofs}

Each identity is stated in the notation of the text, with a pointer to the place where it is used, and proved step by step.
When the text already contains a proof, the proof here is a cleaned-up version of it; when the text only states the identity or
sketches the argument, the proof is new (such proofs carry the comment \texttt{PROOF ADDED BY REWRITER} in the source).
Throughout, ``orthonormal basis'' means a basis $(e_i)$ with $e_i\cdot e_j=0$ for $i\neq j$ and $e_i^2=\varepsilon_i\in\{1,-1\}$
(or $0$ in the degenerate case, where only orthogonality is used).

\paragraph{Used but not proved here.} The following facts are used in the text and are left to the references:
the second Bianchi identity and its contracted form $\nabla^\mu G_{\mu\nu}=0$ \cite{Lee,Carroll}; the Poincar\'e lemma
(Theorem~\ref{thm:dprops}(5)) \cite{Spivak,BottTu}; the existence and uniqueness of $\dd$ (Definition~\ref{def:d}) \cite{Lee};
the normal-coordinate expansion of $g_{ij}$ (Section~\ref{sec:curv}; see the author note there on its sign) \cite{Lee};
and the signs of the Hodge star for indefinite metrics \cite{Carroll}.
